% from ega4-2-fr.tex
\oldpage[IV-2]{79}

\begin{proof}
Il suffit de considérer une extension parfaite $k''$ de $k'$ et en posant
$X''=X\otimes_k k''$, $\mathscr F''=\mathscr F\otimes_k k''$, un point
maximal $x''$ de $\operatorname{Supp}(\mathscr F'')$ qui est au-dessus du
point générique de $Z'$; alors il résulte de (4.7.8) que
$\operatorname{long}(\mathscr F''_{x''})$ est égal aux multiplicités de $Z$ et
de $Z'$.
\end{proof}

\begin{proposition}[4.7.10]
\label{IV.2.4.7.10-fr}
Soient $k$ un corps, $X$ un $k$-préschéma localement de type fini sur $k$,
$\mathscr F$ un $\mathcal O_X$-Module cohérent, $x$ un point de $X$, $z_i$
les points génériques des composantes irréductibles de
$\operatorname{Supp}(\mathscr F)$ qui contiennent $x$. Pour que $\mathscr F$
soit géométriquement réduit au point $x$ (4.6.22), il faut et il suffit que
$x$ n'appartienne à aucun cycle premier associé immergé de $\mathscr F$ et
que, pour tout $i$, la longueur géométrique $\lambda_{z_i}(\mathscr F)$ de
$\mathscr F$ au point $z_i$ soit égale à $1$.
\end{proposition}

\begin{proof}
Montrons d'abord que les conditions de l'énoncé sont suffisantes. En effet, si
$k'$ est une extension finie de $k$, $x'$ un point de
$X'=X\otimes_k k'$ au-dessus de $x$, et
$\mathscr F'=\mathscr F\otimes_k k'$, $x'$ n'appartient à aucun cycle premier
associé immergé de $\mathscr F'$ en vertu de (4.2.7); en outre, en vertu de
(4.7.8), la longueur géométrique de $\mathscr F'$ est égale à $1$ en tout
point générique $z'_j$ d'une composante irréductible de
$\operatorname{Supp}(\mathscr F')$ contenant $x'$, un tel point étant
au-dessus d'un des $z_i$ en vertu de (2.3.4); \emph{a fortiori} on a
$\operatorname{long}(\mathscr F'_{z'_j})=1$, ce qui prouve que $\mathscr F'$
est réduit au point $x'$.

Inversement, il résulte d'abord de (4.2.7) que si, pour une extension finie
$k'$ de $k$, $\mathscr F'$ est réduit en tous les points $x'$ de $X'$
au-dessus de $x$, alors $x$ n'appartient à aucun cycle premier associé immergé
de $\mathscr F$. En outre, en prenant pour $k'$ le corps $k_1$ de l'énoncé de
(4.7.8) (où $x$ est remplacé par un des $z_i$), l'hypothèse que $\mathscr F$
est géométriquement réduit au point $x$ entraîne que
$\lambda_{z_i}(\mathscr F)=1$.
\end{proof}

\begin{proposition}[4.7.11]
\label{IV.2.4.7.11-fr}
Soient $k$ un corps, $X$ un préschéma localement de type fini sur $k$,
$\mathscr F$ un $\mathcal O_X$-Module cohérent, $k'$ une extension de $k$,
$X'=X\otimes_k k'$, $\mathscr F'=\mathscr F\otimes_k k'$, $x$ un point de
$X$, $x'$ un point de $X'$ au-dessus de $x$. Pour que $\mathscr F$ soit
géométriquement réduit (resp. géométriquement ponctuellement intègre) au point
$x$, il faut et il suffit que $\mathscr F'$ soit géométriquement réduit (resp.
géométriquement ponctuellement intègre) au point $x'$.
\end{proposition}

\begin{proof}
Les composantes irréductibles de $\operatorname{Supp}(\mathscr F')$ contenant
$x'$ dominent chacune une composante irréductible de
$\operatorname{Supp}(\mathscr F)$ contenant $x$, et inversement chacune de
ces dernières est dominée par une composante irréductible de
$\operatorname{Supp}(\mathscr F')$ contenant $x'$ ((4.2.7, (i)) et (2.3.4)).
L'assertion concernant la propriété de séparabilité résulte donc de (4.2.7,
(ii)), (4.7.9) et (4.7.10). Soient $Y$ le sous-préschéma fermé de $X$ défini
par l'Idéal $\mathscr J$ annulateur de $\mathscr F$, de sorte que
$Y'=Y\otimes_k k'$ est le sous-préschéma fermé de $X'$ défini par l'Idéal
$\mathscr J'=\mathscr J\otimes_k k'$ annulateur de $\mathscr F'$; dire que
$\mathscr F$ (resp. $\mathscr F'$) est géométriquement ponctuellement intègre
au point $x$ (resp. $x'$) revient à dire que $\mathscr F$ (resp.
$\mathscr F'$) est géométriquement réduit au point $x$ (resp. $x'$) et que
$Y$ (resp. $Y'$) est géométriquement ponctuellement intègre au point $x$
(resp. $x'$). En vertu de l'assertion de l'énoncé concernant la séparabilité,
l'hypothèse que $\mathscr F$ est géométriquement ponctuellement intègre au
point $x$, ou l'hypothèse que $\mathscr F'$ est géométriquement ponctuellement
intègre au point $x'$, entraînent toutes deux qu'il existe dans $Y$ un
voisinage de $x$ qui est séparable, et la conclusion résulte de (4.6.10).
\end{proof}

\begin{definition}[4.7.12]
\label{IV.2.4.7.12-fr}
Soient $k$ un corps, $X$ un $k$-préschéma localement de type fini sur $k$,
$\mathscr F$ un $\mathcal O_X$-Module cohérent, $x$ un point maximal du
support de $\mathscr F$. On appelle \emph{multiplicité}

